\documentclass[12pt]{article}

\usepackage{setspace}
\onehalfspace

%\usepackage[utf8]{inputenc}
%\usepackage{t1enc}
%\usepackage[T1]{fontenc}

\usepackage{amsmath}


%\DeclareMathOperator{\Exp}{Exp}
%\def\Exp{{\mathrm{Exp}}}

%\DeclareMathOperator{\Norm}{{\cal N}}
%\def\Norm{{\cal N}}

%\DeclareMathOperator{\Binom}{Binom}
%\def\Binom{{\mathrm{Binom}}}

%\DeclareMathOperator{\E}{E}
%\def\E{{\mathrm{E}}}

%\DeclareMathOperator{\D}{D}
%\def\D{{\mathrm{D}}}

%\renewcommand{\P}{\operatorname{P}}
%\def\P{{\mathrm{P}}}

%\DeclareMathOperator{\cov}{cov}
%\def\cov{{\mathrm{cov}}}


%\DeclareMathOperator{\corr}{corr}
%\def\corr{{\mathrm{corr}}}


%\renewcommand{\d}{\operatorname{d\! }} %integral \! a nyero
%\def\d{{\mathrm{d\! }}}


%\DeclareMathOperator{\R}{\mathbb{R}}
%\def\R{{\mathbb{R}}}

%\DeclareMathOperator{\Unif}{{\cal U}}
%\def\Unif{{\cal U}}

%\DeclareMathOperator{\Poi}{{Poisson}}
%\def\Poi{{\mathrm{Poisson}}}

%\def\ldots{{...\:}}


\usepackage{moodle}


\begin{document}

\begin{quiz}{gyakorló}

\begin{multi}[feedback={Def.}]{flenvv4}
Ha $\xi$ és $\eta$ diszkrét valségi változóknak létezik a szórása és nemnulla, akkor a 
korrelációjuk
$$
\corr(\xi,\eta)=\frac{\cov(\xi,\eta)}{\D(\xi)\D(\eta)}
$$
\item* igaz
\item hamis
\end{multi}
\begin{multi}[feedback={Ilyesmi a $\corr$-ra igaz.}]{flenvv8}
Ha a $\xi$ és $\eta$ valségi változóknak létezik a kovarianciája, akkor $-1\le \cov(\xi,\eta)\le 1$.
\item igaz
\item* hamis
\end{multi}
\begin{multi}[feedback={Def.}]{flenvv1}
Azt mondjuk, hogy $\xi$ és $\eta$ diszkrét valségi változók függetlenek, ha
$$
\P(\xi=x,\eta=y)=\P(\xi=x)\P(\eta=y)
$$
bármely $x,y\in \R$-re teljesül.
\item* igaz
\item hamis
\end{multi}
\begin{multi}[feedback={Tétel.}]{flenvv3}
Ha $\xi$ és $\eta$ valségi változók függetlenek, akkor
$$
\E(\xi \eta)=\E(\xi)\E(\eta)
$$
feltéve, hogy létezik a szórásuk.
\item* igaz
\item hamis
\end{multi}
\begin{multi}[feedback={Tulajdonság,tétel.}]{flenvv9}
Ha a $\xi$ és $\eta$ valségi változóknak létezik a korrelációja, akkor $-1\le \corr(\xi,\eta)\le 1$.
\item* igaz
\item hamis
\end{multi}
\begin{multi}[feedback={Def.}]{flenvv5}
Ha $\xi$ és $\eta$ diszkrét valségi változóknak létezik a szórása és nemnulla, akkor a 
korrelációjuk
$$
\corr(\xi,\eta)=\E(\xi\eta)-\E(\xi)\E(\eta)
$$
\item igaz
\item* hamis
\end{multi}
\begin{multi}[feedback={Def.}]{flenvv7}
Ha $\xi$ és $\eta$ valségi változóknak létezik a szórása, akkor a 
kovarianciájuk
$$
\cov(\xi,\eta)=\E((\xi-\E(\xi))(\eta-\E(\eta)))
$$
\item* igaz
\item hamis
\end{multi}
\begin{multi}[feedback={Fordítva igaz. Ezek a korrelálatlan (merőleges) változók. }]{flenvv2}
Ha
$$
\E(\xi \eta)=\E(\xi)\E(\eta)
$$
akkor$\xi$ és $\eta$ valségi változók függetlenek.

\item igaz
\item* hamis
\end{multi}
\begin{multi}[feedback={Tulajdonság.}]{flenvv6}
Ha $\xi$ és $\eta$ valségi változóknak létezik a szórása, akkor a 
kovarianciájuk
$$
\cov(\xi,\eta)=\E(\xi\eta)-\E(\xi)\E(\eta)
$$
\item* igaz
\item hamis
\end{multi}
\end{quiz}

\end{document}

